\section{Scheduling DL Inference Workloads}
\label{sec_inference}


\begin{table*}[t]
    \centering
    \caption{\textbf{Summary of schedulers for DL inference workloads in GPU Datacenters.}}
    \vspace{-10pt}
    \resizebox{\textwidth}{!}{

        \begin{tabular}{@{}cccccccccc@{}}
            \toprule
            \textbf{Year}         & \textbf{Scheduler}                           & \textbf{Objectives}            & \textbf{Approaches}                              & \textbf{Advantages}                                    & \textbf{Batching} & \textbf{Colocate} & \textbf{Cloud} & \textbf{\makecell[c]{Exp.             \\ Scale}} & \textbf{\makecell[c]{Source \\ Code}} \\ \midrule
            2017                  & Clipper \cite{Daniel2017clipper}             & \ding{168}\ding{170}\ding{171} & Query-level caching; Layered   architecture      & General abstraction for model selection                & \ding{52}         & -                 & -              & M                         & \ding{52} \\ \midrule
            \multirow{3}{*}{2018} & Space-Time \cite{jain2018dynamic}            & \ding{171}\ding{95}            & GPU sharing across space and time                & Performance isolation under sharing                    & \ding{52}         & \ding{52}         & -              & S                         & -         \\
                                  & Ease.ml   \cite{li2018easeml}                & \ding{168}                     & Multi-tenant model selection                     & Homogeneous declarative   inference platform           & -                 & -                 & -              & -                         & \ding{52} \\
                                  & HiveMind \cite{HiveMind}                     & \ding{171}                     & Sharings of pipelines, weights,   and layers     & Multi-model training and inference                     & \ding{52}         & \ding{52}         & -              & S                         & -         \\ \midrule
            \multirow{8}{*}{2019} & MArk \cite{zhang2019mark}                    & \ding{170}\ding{169}           & Predictive autoscaling on serverless instances   & Flexible to burst requests                             & \ding{52}         & -                 & \ding{52}      & S                         & \ding{52} \\
                                  & Tolerance Tiers   \cite{Halpern2019onesize}  & \ding{168}\ding{170}\ding{169} & Service Version Ensembling                       & Explicit accuracy-latency   trade-off in requests      & -                 & -                 & \ding{52}      & S                         & -         \\
                                  & ParM \cite{ParM}                             & \ding{170}                     & Coded-computation via a learning-based approach  & Erasure-coded resilience for inference                 & \ding{52}         & -                 & -              & M                         & \ding{52} \\
                                  & Gilman et al.   \cite{Challenges19}          & \ding{170}\ding{171}           & Preloading model into GPU memory                 & DNN model execution caching                            & -                 & \ding{52}         & -              & -                         & -         \\
                                  & Nanily \cite{tang2019nanily}                 & \ding{170}\ding{171}           & Adaptive batching; autoscaling                   & Batch size adjustment by remaining time                & \ding{52}         & -                 & -              & S                         & -         \\
                                  & RRL   \cite{qin2019swift}                    & \ding{170}                     & Region-based Reinforcement Learning              & Parallelism configuration tuning                       & \ding{52}         & \ding{52}         & -              & M                         & \ding{52} \\
                                  & TrIMS   \cite{dakkak2019TrIMS}               & \ding{170}\ding{171}\ding{95}  & Multi-layered caching across FaaS                & Memory efficiency by sharing                           & \ding{52}         & \ding{52}         & \ding{52}      & S                         & \ding{52} \\
                                  & Ebird \cite{Ebird}                           & \ding{170}\ding{171}\ding{95}  & CUDA stream paralleism; GPU-side   memory pool   & Transfer-computation overlapping                       & \ding{52}         & \ding{52}         & -              & S                         & \ding{52} \\ \midrule
            \multirow{7}{*}{2020} & GSLICE \cite{dhakal2020GSLICE}               & \ding{171}\ding{95}            & Dynamic GPU resource   apportioning              & Efficient fine-grained sharing                         & \ding{52}         & \ding{52}         & -              & S                         & -         \\
                                  & Clockwork   \cite{Arpan2020clockwork}        & \ding{170}\ding{171}           & Consolidating choice                             & Predictable E2E performance                            & \ding{52}         & -                 & -              & M                         & \ding{52} \\
                                  & Irina \cite{wu2020irina}                     & \ding{170}\ding{171}\ding{95}  & Batching, colocation and preemption              & Graceful general preemption                            & \ding{52}         & \ding{52}         & -              & S                         & -         \\
                                  & PERSEUS   \cite{lemay2020Perseus}            & \ding{170}\ding{169}\ding{171} & Performance and cost characterization on Cloud   & Cost savings under GPU instances                       & \ding{52}         & -                 & \ding{52}      & S                         & \ding{52} \\
                                  & AutoDeep \cite{AutoDeep}                     & \ding{170}\ding{169}\ding{171} & BO and DRL                                       & Cloud configuration and device placement               & -                 & \ding{52}         & \ding{52}      & S                         & -         \\
                                  & DyBatch   \cite{zhang2020dybatch}            & \ding{170}\ding{171}           & Task slicing and reordering                      & Fine-grained batching;   Fairness-driven scheduling    & \ding{52}         & \ding{52}         & -              & S                         & -         \\
                                  & Inferline   \cite{crankshaw2020inferline}    & \ding{170}\ding{169}           & Low-frequency planner; high-frequency tuner      & Near-optimal scaling   cost-efficiency                 & \ding{52}         & -                 & \ding{52}      & M                         & \ding{52} \\ \midrule
            \multirow{5}{*}{2021} & INFaaS \cite{INFaaS}                         & \ding{170}\ding{169}\ding{171} & VM- and model-level autoscaling                  & Automatical model variants selection                   & -                 & \ding{52}         & \ding{52}      & -                         & \ding{52} \\
                                  & Mendoza et al.   \cite{Interference-Aware21} & \ding{170}                     & Latency degradation prediction during colocation & Safe colocation                                        & -                 & \ding{52}         & -              & M                         & -         \\
                                  & Morphling \cite{Morphling}                   & \ding{169}\ding{171}           & Model-agnostic meta-learning                     & Performance prediction under   different configuration & \ding{52}         & \ding{52}         & \ding{52}      & -                         & \ding{52} \\
                                  & Abacus   \cite{cui2021abacus}                & \ding{170}\ding{171}           & Runtime operator scheduling                      & Deterministic latency under   colocation               & -                 & \ding{52}         & -              & M                         & \ding{52} \\
                                  & MIG-SERVING   \cite{tan2021MIG-SERVING}      & \ding{170}\ding{169}           & Greedy algorithm, GA, and MCTS                   & MIG enabled inference scheduling                       & \ding{52}         & \ding{52}         & -              & S                         & -         \\ \midrule
            2022                  & Cocktail \cite{Jashwant2021Cocktail}         & \ding{168}\ding{170}\ding{169} & Weighted majority voting policy                  & Ensembling-based model selection                       & -                 & -                 & \ding{52}      & -                         & -         \\ \bottomrule
        \end{tabular}
    }


    \begin{flushleft}
        \begin{tablenotes}[para,flushleft]
            \scriptsize
            \textbf{Objectives:} \ding{95} Utilization \ding{168} Accuracy \ding{169} Cost \ding{170} Latency  \ding{171} Throughput; \hspace{10pt}  \textbf{Experiment Scale:} S (Single Node) M (Multi Nodes) -: no evaluation on a physical cluster or not clearly specified.
        \end{tablenotes}
    \end{flushleft}
    \vspace{-15pt}
    \label{tab_inference}
\end{table*}

As more DL-based applications are released as online services in our daily life, it becomes more critical to manage and schedule large-scale inference workloads in the GPU datacenter. 
Different from the resource-intensive and long-term training workloads, inference jobs have unique characteristics and requirements (Sec. \ref{sec_background-inference}), which demand new scheduling solutions.
Similar as Sec. \ref{sec_training}, we categorize these inference scheduling techniques based on their objectives, and resource consumption features. Then we give some implications from these works at the end of this section. Table \ref{tab_inference} summaries the relevant papers and their features.





\subsection{Scheduling Objectives}

We first review prior works based on the scheduling objectives.


\subsubsection{Efficiency}
\label{sec_accuracy_latency_cost}


As discussed in Sec. \ref{sec_inference_challenge}, the main objective for scheduling an inference workload is to improve its efficiency. This includes the reduction of inference latency and cost, and improvement of the prediction accuracy {\tiny\encircle{\normalsize I2}}. The challenge here is that there exist trade-offs among different efficiency goals. Here we discuss the techniques to improve each goal as well as to jointly balance and improve them.





\textbf{1) Accuracy efficiency.} Improving the prediction accuracy is a perpetual objective in an inference system. To achieve this, one approach is to collect a set of models, and select the best one to predict the result for each input query.
The scheduling decision made includes model selection and resource allocation among different candidates.
Ease.ml~\cite{li2018easeml} leverages the input and output shape information of the query sample to automatically select the model. It estimates the potential accuracy improvement of each candidate model and then picks the highest one for actual inference. It also formulates the cost-aware model selection process under both single-tenant and multi-tenant settings with multi-armed bandit and Bayesian Optimization. Another effective approach is model ensemble, which combines the prediction results from multiple models to improve the prediction accuracy and generalization. Clipper~\cite{Daniel2017clipper} examines the benefits brought from the model ensemble in computer vision tasks and applies a linear ensemble method to compute a weighted average of the base model predictions. The linear weights are decided by bandit- and learning-based approaches. Rafiki~\cite{wang2018rafiki} leverages an RL model to determine the model set for the ensemble. This model is also used to identify the final optimal model combinations, and tune critical parameters, e.g., batch size.









\textbf{2) Latency efficiency.}
An inference system should have a satisfactory response time, even for burst and fluctuating query requests {\tiny\encircle{\normalsize C7}}.
The latency requirement poses challenges for the scheduler to decide which jobs to be prioritized in the job assignment and rearrangement process.
This objective can be achieved via carefully optimizing the resource allocation.

It is common to launch multiple inference execution instances concurrently to meet the corresponding latency requirement as much as possible due to the low GPU utilization for each request {\tiny\encircle{\normalsize C5}}. Therefore, the inference scheduler can make scheduling decisions aiming at scaling up resources according to the request density to maintain a satisfactory latency.
Clipper~\cite{Daniel2017clipper} conducts linear scaling of inference instances and uses separate docker containers to isolate different models. It replicates the model containers according to the number of queries and applies adaptive batching independently for each model due to the varied execution time.
MArk~\cite{zhang2019mark,zhang2020mark} scales the inference instances with the cloud services. It selects and combines different cloud services like AWS EC2 and Lambda in order based on their prices and scaling abilities. Also, it monitors the system loads and request queuing situations proactively and leverages Lambda to scale up instances when there exist requests violating the latency demands. InferLine~\cite{crankshaw2020inferline} targets the pipelined inference workloads with multiple stages.
It monitors the frequency of queries to each model and makes the scaling decisions of each component separately, to maintain the latency SLOs even during sharp bursts.



A number of works aim to provide bounded latency for inference execution at the system level considering its deterministic execution {\tiny\encircle{\normalsize I1}}. Clockwork~\cite{Arpan2020clockwork} discovers that many DL inference models have deterministic performance because of the underlying deterministic computations. Thus, it guarantees deterministic latency by alleviating the uncertainty introduced by other components of the system. To overcome the uncertainty from memory and cache usages, hardware interactions, and other uncontrollable external performance variations, Clockwork consolidates the configurations among all the system layers during the inference execution, by proactively controlling the memory allocation and deallocation, and disabling concurrent execution of multiple inference workloads to eliminate the interaction.
Reducing the parallelism of execution eliminates the interference from other tasks, but inevitably brings lower throughput and resource utilization. To address this issue, Abacus~\cite{cui2021abacus} tries to guarantee SLO for query requests under the GPU co-location scenarios. It controls the execution sequence and the co-location situation proactively, rather than the default random-ordered execution overlap. Given the explicit order and specific co-location operators on GPUs, Abacus could predict the estimated running time under co-location from the early offline profiling stage. Based on the estimation, the query controller schedules all the simultaneous inference workloads to guarantee the QoS by searching the optimal execution sequence of DNN operators. ParM~\cite{Kosaian2019parity} migrates the concept of erasure codes from distributed computing to model inference systems, and uses learning-based coded computation to introduce redundancy and thus supports the recovery of inference executions with tail latency or failures.

Some solutions proactively schedules the inference workloads and rearranges the execution sequence at the job level. Irina~\cite{wu2020irina} is the first online inference scheduling system, modeling the satisfaction of latency demands as a scheduling problem. By leveraging preemption for DL inference workloads, Irina dynamically decides whether to preempt the ongoing query and launch the later arrived one, which brings significant reduction of average completion time for inference workloads. The main challenge is that existing ML frameworks are not designed and suitable for preemption during execution. Irina carefully manages the preemption process by adding an exit node to the existing dataflow graph of the inference workload, thus enabling safe preemption at arbitrary moments. It is necessary to have more runtime information about the inference workloads for effective scheduling.
Kube-Knots~\cite{Thinakaran2019kubeknots} makes predictions about the resource utilization of each inference workload from two aspects. From the spatial aspect, Kube-Knots discovers the correlations across different resource utilization metrics, and then forecasts the future resource utilization. From the temporal aspect, Kube-Knots predicts the peak inference usage and tries to avoid co-locating jobs which could attain peak consumption of the same resources simultaneously.

\textbf{3) Cost-efficiency.}
The monetary cost becomes one of the main concerns when using public cloud resources to deploy DL inference workloads. Considering the varied compute capabilities and prices for different types of resources and services, a couple of schedulers implement many mechanisms to achieve cost-efficient inference.
MArk~\cite{zhang2019mark,zhang2020mark}
analyzes the cost of utilizing different levels of resource abstractions in Amazon Web Services (AWS) and Google Cloud Platform (GCP) for inference. It finds that the Infrastructure-as-a-Service (IaaS) provides better cost efficiency than Content-as-a-Service (CaaS), while Function-as-a-Service (FaaS) could compensate for the relatively long cold start latency of IaaS at the cost of increased costs. Small instances with advanced CPUs and burstable instances are also recommended. For GPU instances, the cost can be greatly reduced by batch processing as well. Given different levels of capability, scalability, and pricing, MArk greedily selects the most cost-effective type of instances and leverages the spot instances for cost-saving.
AutoDeep~\cite{AutoDeep} considers not only the resource type in the cloud but also the device placement for DL inference. It leverages Bayesian Optimization for the nearly optimal cloud configuration and Deep Reinforcement Learning for the nearly optimal device placement. Constrained by the SLO requirements, AutoDeep performs joint optimization to minimize the total cost of the inference execution. Cocktail~\cite{Jashwant2021Cocktail} develops a distributed weighted auto-scaling policy and leverages the spot instances to minimize the cost.


\textbf{4) Trade-offs between accuracy, latency and cost.}
The objectives of accuracy, latency and cost are not independent {\tiny\encircle{\normalsize C6}}. Improving one goal may possibly compromise another goal if the solution is not designed properly. Besides, users may also have their specific expectations about different objectives. This motivates researchers to explore the trade-offs between these objectives, and devise more flexible and comprehensive scheduling systems.


The adoption of multiple models can improve the model inference accuracy, but might also increase the response latency and cost. Several works track the latency and prediction accuracy of different models and implement mechanisms to select the most appropriate ones determined by the schedulers.
Clipper~\cite{Daniel2017clipper} introduces a model selection abstraction, which supports both single model selection and model ensemble selection. It executes the inference for all the models and combines their results. It observes the corresponding accuracy and latency feedback continuously to make the selection with a best-effort search method. Model-Switching~\cite{jeff2020modelswitching} pursues the trade-offs between computational cost and service accuracy by switching different model variants proactively, to improve the accuracy of responses under the latency constraint. By maximizing the ratio of correct predictions returned within the deadline, it makes selections among model variations with different computation demands and accuracy.  Cocktail~\cite{Jashwant2021Cocktail} balances the cost with accuracy and latency on the public cloud via the optimization of the model ensemble. With a dynamic model selection policy that searches models tightly with the accuracy and latency bounds, Cocktail reduces the candidates in the ensemble and accomplishes fast and efficient model selection.







Some schedulers allow users to specify their demands about accuracy, latency and cost, and make scheduling decisions directly according to the demands. Tolerance Tiers~\cite{Halpern2019onesize} discloses the efforts the system can offer to achieve each objective, and makes users programmatically select their demands. Observing that improving the accuracy of some extreme requests can increase the latency greatly, Tolerance Tiers relaxes and sacrifices the accuracy demand to improve the latency and service cost. Each tier defines an error tolerance to indicate the tolerable accuracy loss, and an optimization objective. Then, Tolerance Tiers optimizes the objective under the constraint of the maximum error tolerance. INFaaS~\cite{Neeraja2019INFaaSw,INFaaS} also asks for the performance and accuracy demands from the users.
It generates some variants from the existing models with different specific parameters (e.g., batch size, hardware configurations,  hardware-specific parameters). After one-time profiling for each variant, INFaaS selects the model variant based on the resource consumption and performance information from the profiling to serve users' requests. Since each model variant may have different performance and monetary costs during execution, INFaaS makes the selection via a heuristic-based approach, which selects the variant with the minimum cost while meeting the SLO constraint, or upgrades existing variants with higher performance to fulfill the burst queries.


\subsubsection{System Throughput}
\label{sec_throughput}

Another important objective for scheduling inference workloads is to improve its throughput capability. The techniques to achieve this goal is summarized as follows. 



\textbf{1) Batching execution.}
One common approach is to batch multiple inference queries, and execute them concurrently. Handling individual inference queries usually leads to GPU underutilization {\tiny\encircle{\normalsize C5}}. Hence, batching inference can efficiently improve the utilization and reduce the system overhead. Like job queuing in parallel job scheduling, batching multiple queries can delay the execution of the requests that come earlier, and possibly jeopardize the SLO requirement. Setting a proper batch size is critical to balance such delays and system throughput. Most schedulers dynamically adjust this hyperparameter based on the actual SLO requirement and queuing situation.





First, some schedulers adopt heuristic methods to tune the batch size. Clipper~\cite{Daniel2017clipper} and Rafiki~\cite{wang2018rafiki} apply the practical Additive-Increase-Multiplicative-Decrease (AIMD) algorithm to adjust the inference batch size. Specifically, the batch size is additively increased by a fixed amount until the latency of processing a batch exceeds the latency requirement and then multiplicatively decreased by a fixed percent. Clipper evaluates that AIMD is simple yet effective and adaptive to the changing throughput of a model in special cases. It also aggressively delays the execution of queries under moderate loads to the subsequent batch, which can bring a significant throughput increase for some models. GSLICE~\cite{dhakal2020GSLICE} also applies a similar adaptive and self-learning approach to determine the optimal batch size. It carefully tracks the execution time and increases the batch size until the execution of the last batch exceeds the SLO requirement. Ebird~\cite{Ebird,E2bird} proposes a novel elastic batching mechanism, which runs different CUDA streams concurrently for different batches. Motivated by the observation that processing multiple queries in a large batch has similar performance as multiple CUDA streams with smaller batch sizes, Ebird dynamically adjusts the batch size per stream to utilize the spare GPU resources and fulfill the whole GPU. In each scheduling round, it selects and launches a job based on its batch size and remaining GPU resources in a best-effort manner, squeezing the full GPU utilization and minimizing inference queuing delay.

Second, some schedulers propose optimization-based methods to balance the inference delay and throughput.
MArk~\cite{zhang2019mark,zhang2020mark} considers the maximum time of delaying a query, profiles the processing rate without batching, and searches for the optimal batch size under the SLO constraint. Nanily~\cite{tang2019nanily} presents the upper bound of the batch size by retrieving the maximum remaining time for the requests, calculated as the remaining time towards the deadline subtracted by the least queuing time for the available resources. It then derives the corresponding batch size, which makes the inference execution time equal or close to the maximum remaining time. DyBatch~\cite{zhang2020dybatch} considers the fairness of the delay for each independent workload when batching. It implements fine-grained batching schemes along with fairness-driven scheduling, which can compensate for the deviation of slowdown for small inference workloads. DyBatch organizes the workload batches in a time-sharing manner and selects the batch with the lowest resource utilization for running, thus maintaining fairness and minimizing the slowdown of each workload.

\textbf{2) Caching and reusing. } Another widely-used strategy for throughput improvement is caching and reusing the prediction results across different requests {\tiny\encircle{\normalsize C7}}.
The scheduler selects the request that benefits most from caching and allocates proper resources. This can be done at two levels.

The first direction is to perform optimization at the query level. To provide fast responses to different queries, the inference system can cache the inference execution and prediction results for burst queries. Clipper~\cite{Daniel2017clipper} maintains a prediction cache for requests with the same target model and the query input. Then it can produce the results for some queries without evaluating the model, thus increasing the inference throughput. Clipper also applies an LRU cache eviction policy to optimize the caching efficiency. However, this approach may be less efficient when the queries do not have high similarities in practical scenarios, which leads to high cache miss rates and evictions.


The second direction is to perform optimization at the device level.  Gillman \etal~\cite{Challenges19} proposed to cache the DL models instead of the inference results. It schedules models to be loaded into the limited GPU memory to maximize the probability of servicing an incoming request without swapping the models in and out of the memory, thus accelerating the inference by eliminating the cold start latency with cache hits. The caching and eviction policy considers many runtime aspects of DL inference workloads, including model size, frequency, model accuracy, and speed. This work also discusses some future directions for more dynamic caching mechanisms and policies, like framework-level GPU memory-friendly optimization, proactively loading and evicting, and cluster-level GPU memory allocation. To address the limitation of GPU memory, GSLICE~\cite{dhakal2020GSLICE} and HiveMind~\cite{narayanan2018accelerating} explore the common GPU memory component in different inference models, and propose to save GPU resources via memory sharing. Particularly, GSLICE enables efficient GPU memory sharing by allowing the reuse of model parameters via modifications to the DL framework, exposing the CUDA address to different instances. Therefore, it supports loading the inference model-related parameters only once to the GPUs, resulting in faster module loading. HiveMind extends the shared content and brings more possibilities of shared model weights and layers across different inference workloads, saving the overhead from both model loading and inference evaluation. TrIMS~\cite{dakkak2019TrIMS} organizes the memory sharing of different models in a more systematic design. The model resource manager in TrIMS offers a multi-tiered cache for DL models to be shared across users' FaaS functions and enables DL model sharing across all levels of the memory hierarchy in the GPU, CPU, local storage, and remote storage in the cloud. TrIMS reconciles the lifecycle of model memory consumption and carefully handles the cache misses and evictions. It also considers multi-node, isolation, and fairness problems during sharing. Extensive evaluations on different models show its general abilities to improve the inference performance by mitigating the model loading overhead.




\textbf{3) System configuration tuning. }
Besides the optimization techniques detailed above, there exist some schedulers leveraging end-to-end configuration tuning to improve the system throughput. Morphling~\cite{wang2021Morphling} formulates the optimal configuration search as a few-shot learning problem. Then it adopts model-agnostic meta-learning (MAML)~\cite{finn2017model} to combine offline meta-model training for inference serving performance modeling under varied hardware and runtime configurations, and performs online few-shot learning to predict the service performance. Based on the prediction, Morphling auto-tunes the resource provisioning configurations and makes better scheduling decisions.
RRL~\cite{qin2019swift} concentrates on optimizing the parallelism configurations from different levels, including request level parallelism and intra-request level (inter-op and intra-op) parallelism, which have strong impacts on the latency of the entire system. RRL utilizes a region-based RL method to tune the parallelism configurations and reduce the inference processing latency, based on the system performance similarity between different configurations within a similar parallelism setting.




\subsection{Resource Consumption Feature}
\label{sec_resource}
Similar to training workloads, the inference schedulers can also be categorized based on the resource consumption feature. Below we detail the scheduling solutions designed to target these features.


\subsubsection{Colocation and resource sharing}
\label{sec:inference_colocation}

From Challenge {\tiny\encircle{\normalsize C5}} in Sec. \ref{sec_inference_challenge}, executing one inference request can lead to GPU resource underutilization. The recent development of GPU architecture designs motivates the GPU sharing from both the hardware perspective~\cite{mps,nvidia2020a100whitepaper} and software perspective~\cite{shi2012vCUDA,Salus}. GPU sharing across different inference requests can significantly improve the resource utilization by better leveraging GPUs' parallel compute capability. However, it can also incur the risks of violating the latency requirements due to the uncertain running time.

One line of research works adopt the static colocation strategy to guarantee the inference latency. Space-Time~\cite{jain2018dynamic} calls for both space-sharing and time-sharing
for DL inference with GPUs. It preserves the predictability and isolation during virtualization by monitoring the inference latency per kernel. Then it improves the utilization by merging concurrent small kernels into larger super-kernels that can fill up the GPU utilization under the time-sharing mechanism. Irina~\cite{wu2020irina} only considers a safe GPU colocation situation based on the peak GPU requirement of the workloads, when their total GPU requirement does not exceed the capacity. It assumes there is no interference and slowdown on the job completion time and heuristically places the newly arrived job on the GPU with the smallest JCT. Nexus~\cite{shen2019nexus} also applies a heuristic approach to select the requests to be co-located on the same GPU. First, it identifies the most appropriate batch size for throughput and SLO needs of the existing inference workload. Afterward, it establishes all possible combinations within a GPU’s duty cycle on a single GPU in a best-fit manner, and maximizes the utilization without violating the latency requirement.



Some other works introduce dynamic colocation mechanisms for managing the inference workloads.
GSLICE~\cite{dhakal2020GSLICE} is an inference system to systematically achieve safe and efficient GPU sharing. It leverages spatial GPU multiplexing for different inference requests on top of the state-of-the-art GPU spatial multiplexing framework MPS. It intensively evaluates the colocation performance with and without MPS, and with the resource provisioning limits. It discovers that the colocation interference could be amortized under the careful configuration of the resource limits. The performance improvement reaches a point of diminishing returns (i.e., kneepoint) after certain configurations, which has a non-linear relationship with the throughput and latency of the inference. Based on these observations, GSLICE tracks the kneepoint of different inference workloads and partition the whole GPU according to their kneepoints. It also designs a hot-standby mechanism to dynamically adjust the resource limit configuration of the specific inference job, along with other implementation optimizations including minimizing the data transfer overhead. MIG-SERVING~\cite{tan2021MIG-SERVING} leverages the hardware support for GPU virtualization (i.e., MIG \cite{mig} on NVIDIA A100\cite{nvidia2020a100whitepaper}) for efficient GPU colocation.
With MIG, an A100 card could be dynamically partitioned into several instances with smaller hardware capacities under some hard constraints. MIG-SERVING discovers that the throughput of most models does not grow linearly with the increase of resources on different instances.
It establishes a reconfigurable scheduling problem and applies a generic algorithm to find a sub-optimal and feasible solution, and improve it via a search-based method.


Since the colocation of multiple inference workloads multiplexes the GPU, it is hard to measure and predict their running time due to the colocation interference.
Therefore, several inference scheduling systems focus on estimating and handling the colocation interference. INFaaS~\cite{Neeraja2019INFaaSw,INFaaS} targets the inference services in the cloud. It identifies the colocation interference caused by the shared hardware resources.
Then it allocates the available resources to the interfered instances by migration or VM-level scaling. The evaluation shows that INFaaS can save the monetary cost by GPU sharing and satisfy the latency requirements by VM-level scaling. Mendoza \etal~\cite{Interference-Aware21} proposed an interference-aware scheduler for inference workloads, which proactively considers the impact of interference on the latency from co-location. By predicting the performance degradation, it minimizes the latency degradation and reduces the SLO violations. It improves the prediction accuracy by exploiting the similarity of co-location configurations, including inference model attributes and machine types. PERSEUS~\cite{lemay2020Perseus} compares the cost and throughput under exclusive execution and colocation for inference workloads. It concludes that the mixed ratio of different models affects the cost efficiency of colocation. The interference on the model and data loading time during cold start also differs across different models because of the cache and data transfer requirements under colocation.






\subsubsection{Heterogeneous Resources.}
Several works exploit both CPU and GPU machines of different sizes for DL inference, especially for the cloud environment. MArk \cite{zhang2019mark,zhang2020mark} characterizes and compares the cost and performance of inference workloads on different types of instances available in GCP and AWS. It concludes that smaller instances with advanced CPU models achieve a higher performance-cost ratio. It also suggests that GPU instances can achieve lower per-request cost and smaller inference latency than CPU instances with appropriate batch sizes. PERSEUS \cite{lemay2020Perseus} performs a similar cost-efficiency characterization, considering instances with multiple onboard GPUs. It examines that some DL models with high GPU utilization could introduce intensive interference with other inference workloads in multi-GPU instances. AutoDeep \cite{AutoDeep} and Cocktail \cite{Jashwant2021Cocktail} focus more on the price of different cloud configurations and search for the best configuration according to the pricing information.






\subsection{Implications}


In Sec. \ref{sec_accuracy_latency_cost} we mainly discuss the monetary cost of deploying DL inference workloads. In a GPU datacenter, how to improve the efficiency of energy cost is also critical. While this objective has been extensively explored for training workloads (Sec. \ref{sec_training_efficiency}), it is relatively less studied for the inference workloads. Some works \cite{facebook,FaceBookTrace} provide some energy characterizations of production DL inference clusters. Kube-Knots~\cite{Thinakaran2019kubeknots} presents simple energy efficiency comparisons of inference workloads between GPUs and CPUs. It is necessary to comprehensively explore the energy optimization of different DL inference models with different types of compute resources, and design more sophisticated energy-saving mechanisms with the consideration of latency and resource utilization. This will be a promising research direction in the future.





Most of the above works treat the inference jobs as a black box for management and optimization. In reality, an inference pipeline may consist of several separate stages to fulfill the query. It is interesting to consider the characteristics of internal stages to optimize the execution in a white-box manner. PRETZEL~\cite{PRETZEL} leverages this idea, which stores and re-uses the model parameters and computation among similar white-box representations of pipeline stages to reduce the resource utilization. We expect more future works will focus on the optimization of inference workloads at this sub-request level.




As the scale of DL models grows faster than the GPU compute capacity, it is more difficult to accommodate single models on a single GPU or machine. This problem becomes more serious when cloud users adopt the resource-constraint serverless platform for inference services. A natural solution to this problem is model partitioning,
Gillis~\cite{yu2021Gillis} splits the network layers and minimizes the inference latency via dynamic programming. It also adopts some searching methods like Bayesian Optimization and RL to minimize the cost. AMPS-Inf~\cite{Jananie2021AMPS-Inf} jointly considers the model partitioning and resource allocation by calculating the optimal execution and resource provisioning plans under the constraint of the response time in the serverless platforms. It is worth more effort to explore the methodologies of serving large models with limited resources for different objectives.





Some inference systems for CPU clusters predict the future trends of query requests to satisfy the latency requirement. For instance, BARISTA~\cite{Bhattacharjee2019BARISTA} predicts the future workload patterns based on the historical data and estimates the required resources for the application to maintain the SLOs. Then it makes resource provisions based on the difference between the desired throughput and latency requirement with the current ones. Swayam~\cite{gujarati2017Swayam} ensures an appropriate replica of service instances by predicting the load and auto-scales the service in a fully distributed way along with the self-reclamation of resources. It adopts linear regression to predict the request arrival rate over short periods to tolerate the rapid changes in time. Since these solutions are hardware-agnostic, they can be applied to the GPU clusters as well. We also expect this strategy can inspire more solutions dedicated to the GPU clusters.


